\section{Kato vanishing and the actual type-II rigidity producer}
\label{sec:type-ii-rigidity-v47}

We prove the original \texttt{KatoTwo} and \texttt{RigII} assertions.
The latter concerns the full nine-dimensional polarized adjoint of an
actual rank-three deformation point.  Its modularity hypothesis is
retained exactly: every rank-two extension specified by
\texttt{Dat.IsExt} is modular for the supplied field isomorphism
\(\iota:\overline{\mathbf Q}_\ell\simeq\mathbf C\).
We construct one such extension and prove its required properties.
In particular, no adjoint-vanishing, newform-to-base-change, or Kato
predicate is left as an unnamed producer.

All source references are at commit
\texttt{7cbf7a0c32fca1badf05931879a37c3f8a2b03b4}.
The original tied period-ring model, finite Galois cohomology,
class-field theory, and the explicitly cited published automorphic
theorems remain classical foundations.  The proofs below do not
reconstruct those foundational theorems.  They reconstruct their
applications, coefficient comparisons, and every finite-dimensional
consumer needed for these two endpoints.

Write \(G=G_{\mathbf Q}\), \(H=G_K\), and
\(\epsilon:G\to\mathbf Z_\ell^\times\) for the cyclotomic
character.  Frobenius in an attachment formula is geometric, so
\(\epsilon(F_p)=p^{-1}\).  At a fixed point let
\(\eta:G\to\{\pm1\}\) be the quadratic character of the
original imaginary quadratic field \(K\), and let \(c\) be the
specified complex conjugation.  The original coefficient ring is
\(\mathcal O\); it is not replaced when constructing the point.
We use a finite extension of \(\mathbf Q_\ell\) only to apply
classical rational cohomology theorems, and return each primitive
to the original algebraic closure.

\subsection{The precise Selmer and Kato comparison}

% BEGIN KTR1
\paragraph{KTR1. Finite coefficients, functorial local conditions, and finite discrete Selmer groups.}
\label{claim:KTR1}
The source's \texttt{BlochKatoSelmerVanishes} is rational
Bloch--Kato vanishing with its full original cocycle scope.  The
following three operations preserve that scope: a continuous
equivariant linear map, finite coefficient extension, and restriction
to a finite number-field extension.  Moreover, if \(V\) is a
finite-coefficient representation and the discrete Bloch--Kato
Selmer group attached to a stable lattice \(T\subset V\) is
finite, then \(H^1_f(L,V)=0\).

\emph{Proof.}
ADK1 puts a continuous algebraic-closure-valued representation and
any one continuous cocycle over one finite coefficient field.
This uses compactness and Baire's theorem, not an assertion that
an arbitrary compact subset of the algebraic closure is contained
in a finite extension.  ADK5 identifies the source's chosen
inertias and decomposition-group graphs with the classical groups.
The graph identification is continuous because the inverse image
of the stabilizer of a finite set of global algebraic elements
fixes a finite local extension.  With BX1's tied period model, the
local condition at \(\ell\) is exactly the kernel of
\[
 H^1(L_v,V)\longrightarrow
 H^1(L_v,V\otimes_{\mathbf Q_\ell}B_{\rm cris}),
 \tag{KTR.1}
\]
and away from \(\ell\) it is the kernel of restriction to inertia.
Positive cohomology at real places vanishes by averaging.

If \(f:V\to V'\) intertwines the actions, apply \(f\) to a
cocycle and to each inertia primitive.  At \(\ell\), apply
\(f\otimes1\) to a crystalline primitive.  On a pure tensor,
\[
 (f\otimes1)(g v\otimes g b)=g f(v)\otimes g b,
\]
so this map commutes with the diagonal action, and the identity
extends by linearity.  Restriction of a primitive is again a
primitive, proving the restriction assertion.  Local extension of
scalars causes no hidden enlargement of the kernel: every given
algebraic-closure primitive uses finitely many vector coefficients,
contained in some finite \(E'/E\).  An \(E\)-linear retraction
\(E'\to E\) fixing \(1\), tensored with \(1\) on
\(B_{\rm cris}\), brings that primitive back to \(E\).
Different places may use different \(E'\).

For the last assertion choose a finite set \(S\) containing the
ramified places, the places above \(\ell\), and the infinite
places.  A rational Selmer class factors through \(G_{L,S}\):
outside \(S\) the action is trivial on inertia, so its
unramified condition says that its restriction to inertia is
zero, and the cocycle descends through their closed normal closure.
The finite-coefficient global cohomology foundation in DX1 gives
that \(H^1(G_{L,S},T)\) is a finitely generated
\(\mathbf Z_\ell\)-module.  Equivalently this follows from
finite cohomology of \(T/\ell T\), inverse limits, and
topological Nakayama.  Its image in \(H^1(G_{L,S},V)\) is a
bounded lattice.  The map to \(H^1(G_{L,S},V/T)\) sends the
rational Selmer group into the discrete Selmer group whose local
conditions are the images of the rational local conditions.
Its kernel is contained in this bounded lattice, by the long
exact sequence for \(0\to T\to V\to V/T\to0\).
The image of a divisible group in a finite group is zero.
Thus the whole rational vector space \(H^1_f(L,V)\) is contained
in a bounded lattice.  A nonzero vector and all its
\(\ell^{-m}\)-multiples cannot be so contained.  The vector
space is zero, as asserted.  This argument does not require the
order of the finite discrete Selmer group to be prime to \(\ell\).
% END KTR1

% BEGIN KTR2
\paragraph{KTR2. The exact noncentral critical Kato application.}
\label{claim:KTR2}
Let \(e\ge1\), let \(f\) be an original normalized newform
of weight \(k=2e+1\), and suppose the continuous representation
\(W\) satisfies
\[
             f.\mathrm{Att}\bigl(\iota,W\epsilon^{-e}\bigr).
 \tag{KTR.2}
\]
For every finite abelian extension \(L/\mathbf Q\),
\(H^1_f(L,W)=0\), with the exact coefficient and local meaning
of KTR1.

\emph{Published theorem and normalization check.}
Kato, \emph{p-adic Hodge theory and values of zeta functions of
modular forms}, Ast\'erisque 295 (2004), Theorem 14.2(1),
states finiteness of \(\operatorname{Sel}(L,T)\) for a stable
lattice in \(V_{F_\lambda}(f)(r)\), when \(L/\mathbf Q\)
is finite abelian, \(1\le r\le k-1\), and \(2r\ne k\).
It applies to every finite place \(\lambda\) of the eigenvalue
field, with no ordinarity assumption and no lower bound on its
residue characteristic.  The paper includes the CM case; the
separate treatment of that case does not remove it from the
statement.  The local Selmer convention is the one in \S14.1:
the local conditions on the discrete module are images of the
unramified rational conditions away from \(\ell\) and of the
\(B_{\rm cris}\) kernel at \(\ell\).
The primary text is
\url{https://www.numdam.org/item/AST_2004__295__117_0/}.

The original good-prime eigenrelations, normalization
\(a_1=1\), nebentype transformation, and exclusion of every
smaller conductor are the newform data used in RNC2's
adelization.  Strong multiplicity one fixes the cuspidal
eigensystem; the conductor condition selects its one-dimensional
newvector line.  Thus this is the classical normalized newform
to which Kato applies, with its finite eigenvalue field, including
the cases where its level is divisible by \(\ell\).

The cohomological realization \(V_{F_\lambda}(f)\) has geometric
Frobenius polynomial, by the explicit inverse-arithmetic-Frobenius
formula (14.10.3) on printed page 240,
\(X^2-a_p(f)X+\chi_f(p)p^{k-1}\).
This is exactly the polynomial in the source's
\texttt{Newform.heckePoly}; it is not its reciprocal polynomial.
The source \texttt{Att} includes semisimplicity, unramifiedness
at all but finitely many primes, and this same polynomial for
every prolongation and every arithmetic Frobenius inverted.
The usual newform Galois realization and (KTR.2), after finite
coefficient descent, therefore agree by Chebotarev and the
semisimple character criterion (CAP6).  Hence
\[
          W\simeq V_{F_\lambda}(f)(e).
 \tag{KTR.3}
\]
The chosen \(\lambda\) is the one induced by the fixed
\(\iota^{-1}\) on the eigenvalue field.  If the actual
coefficient field is larger than \(F_\lambda\), restriction
of scalars is a finite direct sum, so the same vanishing follows
after extension.  There is no change of \(\iota\).
Now \(1\le e\le2e=k-1\) and \(2e\ne2e+1=k\).
Kato's theorem and KTR1 prove (KTR.3)'s rational Selmer vanishing.
This is an application at \(r=e\), rather than at the excluded
central value.  Its proof supplies the actual coboundary for
each original continuous cocycle.
% END KTR2

% BEGIN KTR3
\paragraph{KTR3. Finite-character descent and the full original KatoTwo.}
\label{claim:KTR3}
Under KTR2, for every continuous finite-order abelian character
\(\varphi:G\to\overline{\mathbf Q}_\ell^\times\),
\[
                       H^1_f(\mathbf Q,W\varphi)=0.
 \tag{KTR.4}
\]
Consequently the full original \texttt{KatoUse} and
\texttt{KatoTwo} hold, with every coefficient closure, embedding,
level, weight, and allowed coefficient prime in their statements.

\emph{Proof.}
Let \(L\) be the finite abelian extension fixed by
\(\ker\varphi\).  On \(G_L\) the twist is \(W\), so
restriction of a Selmer cocycle is a coboundary by KTR2 and
KTR1.  Subtract that coboundary.  The remaining cocycle vanishes
on \(G_L\); its values are invariant under \(G_L\), and it
descends to \(\operatorname{Gal}(L/\mathbf Q)\).
Every cocycle of a finite group with characteristic-zero vector
coefficients is a coboundary.  Explicitly, for a cocycle \(a\)
on a finite group \(Q\), put
\(v=-|Q|^{-1}\sum_{q\in Q}a(q)\); then
\(a(g)=g v-v\).  Thus restriction is injective in degree one.
No condition \(\ell\nmid[L:\mathbf Q]\) is needed.

For \texttt{KatoUse}, the original condition
\texttt{KillsMu}\((M,\varphi)\) makes \(\varphi\) factor
through \(\operatorname{Gal}(\mathbf Q(\mu_M)/\mathbf Q)\),
a finite abelian quotient.  It is continuous and satisfies the
condition just proved.  The original assumption \(M\ge1\)
and the original odd-prime quantifier are left intact.

For \texttt{KatoTwo}, retain all fields of the source's
\texttt{Used} record: \(\ell\ne2\), \(e\ge1\),
\(\ell-1\nmid2e\), imaginary quadratic \(K\),
\(\ell\nmid|d_K|h_K\), the even positive level and its
coprimality, both character conditions, the complex-conjugation
condition, the exact \texttt{Level} and \texttt{OrdPoly},
determinant, inertial trace, residual trace, the specified flag,
and irreducibility on the infinite cyclotomic kernel.
The \texttt{Level} field includes continuity of \(W\).
For this conclusion the extra hypotheses are unnecessary:
the character \(\theta\) has kernel \(G_K\), hence is the
continuous quadratic character, and we may apply (KTR.4) to
\(\varphi=1\) and \(\varphi=\theta\).
This gives the source's two literal column representations
\(\operatorname{colRep}(\operatorname{tw}W1)\) and
\(\operatorname{colRep}(\operatorname{tw}W\theta)\).
The claim is the complete implication with the original
\texttt{Used} antecedent, not a demand for a stronger antecedent.
It includes \(\ell=3\) wherever the original antecedent is
satisfiable; the published theorem itself excludes no such prime.
% END KTR3

\subsection{The actual point and the smaller local calculation}

For KTR4--KTR13 fix all the original data in \texttt{RigII}:
\(D:\texttt{Dat}(\Omega,K,\mathcal O,e,\psi)\),
\(U:D.\texttt{ULoc}\), the marked element \(g_1\), and
an actual \(r\) satisfying \(D.\texttt{IsPt}(U,g_1,r)\).
Assume \(\texttt{TypeII}(\eta,r)\).  Let \(R(g)\) be the
matrix part of \(r(g)\).  On \(H=\ker\eta\) it is an
ordinary representation; its only possible off-diagonal blocks
are the corner positions \(0,2\).  Write
\[
 R(h)=\begin{pmatrix}a_h&0&b_h\\0&\alpha_h&0\\c_h&0&d_h\end{pmatrix},
 \qquad B_h=\begin{pmatrix}a_h&b_h\\c_h&d_h\end{pmatrix}
 \quad(h\in H).
 \tag{KTR.5}
\]
The corner representation is irreducible.  The point is continuous,
its entries are integral and entrywise congruent to the original
\(r_0(\eta,\Psi)\), and \(R(c)=1\).
These are fields of \texttt{IsPt} and \texttt{Normalised}.
Its away-from-\(\ell\) killing condition and its map through
the original local ideal \(I_t\) will both be used.

% BEGIN KTR4
\paragraph{KTR4. The original local ideal forces an ordered inertia flag on a finite-index subgroup.}
\label{claim:KTR4}
Choose the \(\ell\)-adic embedding stored in \(D\), let
\(I\subset G\) be its inertia group, let \(\tau\) be the
stored finite local character, and put \(J=I\cap\ker\tau\).
Then \(I\subset H\), \(\tau^2|_I=1\), and on \(J\)
the actual \(R\) has an ordered flag with characters
\[
                      (\epsilon^e,1,\epsilon^{-e}).
 \tag{KTR.6}
\]
Every element of \(I\) preserves this ordered flag.

\emph{Proof.}
The assumptions \(\ell\nmid d_K\) and \(K/\mathbf Q\)
quadratic put inertia in \(H\).  The original \texttt{GD}
field \((\operatorname{comp}_\ell\psi)^2=1\) gives the
assertion about \(\tau\), and the stored local character
identity is \(\Psi|_I=\epsilon^e\tau\).
The embedding comparison in \texttt{DatLoc} follows at finite
Galois levels by extending the specified \(\ell\)-adic
valuation and taking its decomposition group; inverse limits
identify the stored local homomorphism, inertia, and cyclotomic
character with these groups.  This uses the actual embedding
equation \(s(jx)=j(\iota_s x)\) of \(D\), so it does not
choose a different local action.

Unfold the defining ideals.  The ordinary ideal is contained in
its saturation \(I_Y\) (use exponent zero), which is contained
in its radical \(I_{\rm cr}\).  Every minimal prime used in
\(I_t\) contains \(I_{\rm cr}\), so
\[
                    I_{\rm ord}\subset I_t.
 \tag{KTR.7}
\]
This also holds if that set of primes is empty.  The local map
in \texttt{IsPt} kills \(I_t\); consequently for all
\(s,s_0,s_1,s_2\in J\) it gives
\[
 \det(T-R(s))=\prod_{a=0}^2(T-\epsilon(s)^{e(1-a)}),
 \qquad
 \prod_{a=0}^2\bigl(R(s_a)-\epsilon(s_a)^{e(1-a)}1\bigr)=0.
 \tag{KTR.8}
\]
The order of the three factors is as displayed.

The image of \(\epsilon\) on rational \(\ell\)-inertia is
all \(\mathbf Z_\ell^\times\).  Indeed each cyclotomic
\(\ell\)-power extension of \(\mathbf Q_\ell\) is totally
ramified, and compactness passes the finite surjections to the
inverse limit.  Choose \(u\in I\) with
\(\epsilon(u)=1+\ell\) and replace it by \(s_*=u^2\).
Then \(s_*\in J\) and \(z=\epsilon(s_*)^e\) has infinite
order.  The three eigenvalues \(z,1,z^{-1}\) are distinct.
Diagonalize \(R(s_*)\) in precisely that order, writing
\(D_* =\operatorname{diag}(z,1,z^{-1})\).

In that basis insert \((s,s_*,s_*)\),
\((s_*,s,s_*)\), and \((s_*,s_*,s)\) in the second
equation of (KTR.8).  If \(N=R(s)\) and
\(p=(\epsilon(s)^e,1,\epsilon(s)^{-e})\), the equations are
\[
 \begin{split}
 &(N-p_0)(D_*-1)(D_*-z^{-1})=0,\\
 &(D_*-z)(N-p_1)(D_*-z^{-1})=0,\\
 &(D_*-z)(D_*-1)(N-p_2)=0.
 \end{split}
 \tag{KTR.9}
\]
Here a scalar denotes its scalar matrix.  The first equation
gives \(N_{00}=p_0,N_{10}=N_{20}=0\); the second gives
\(N_{11}=p_1\); the third gives
\(N_{21}=0,N_{22}=p_2\).  The factors divided out are
differences among \(z,1,z^{-1}\), all nonzero.  Thus the
matrices on \(J\) are upper triangular with (KTR.6).

Finally \(J\) is normal in \(I\) and \(\epsilon\) is
unchanged by conjugation.  For \(g\in I\), the matrices of
\(s_*\) and \(g s_*g^{-1}\) are upper triangular with the
same distinct ordered diagonal.  Their intertwiner \(R(g)\)
is upper triangular: the \((2,0)\) equation first kills that
entry, then the \((2,1)\) and \((1,0)\) equations kill those
entries.  Their coefficients are again nonzero diagonal
differences.  This proves preservation of the flag by \(I\).
Only (KTR.7), rather than the stronger full local crystalline
component theorem, has been required.
% END KTR4

% BEGIN KTR5
\paragraph{KTR5. The actual middle and determinant characters are trivial on the quadratic field.}
\label{claim:KTR5}
For every \(h\in H\), the matrix in (KTR.5) satisfies
\[
                          \alpha_h=1,\qquad \det B_h=1.
 \tag{KTR.10}
\]

\emph{Proof.}
Both \(\alpha\) and \(\delta=\det R|_H\) are continuous
characters.  Integral entrywise residual closeness implies
\(\alpha(h)\equiv1\) and \(\delta(h)\equiv1\) modulo
the maximal ideal for every \(h\in H\): the residual diagonal
on \(H\) is \((\overline\Psi,1,\overline\Psi^{-1})\).
We shall prove both characters unramified everywhere over \(K\).

First consider the inertia group in KTR4.  The condition
\(\ell-1\nmid2e\) supplies a unit
\(a\in\mathbf Z_\ell^\times\) with
\(\bar a^{2e}\ne1\).  Choose \(s\in I\) with
\(\epsilon(s)=a\) and put \(t=s^2\in J\).
Then \(\overline{\epsilon(t)^e}\ne1\).
For any \(x\in J\), the middle-block eigenvalue
\(\alpha(x)\) is one of the roots in (KTR.8).
If \(\overline{\epsilon(x)^e}\ne1\), residual closeness
forces \(\alpha(x)=1\).  In the other case apply this
argument to \(xt\) and to \(t\), and divide their character
values.  Hence \(\alpha|_J=1\).  Also (KTR.8) gives
\(\delta|_J=1\) directly.  For \(x\in I\), its square
lies in \(J\).  Each of \(\alpha(x),\delta(x)\) is thus
a square root of unity congruent to \(1\), and \(\ell\ne2\)
makes it equal to \(1\).  The other places above \(\ell\)
are conjugate in \(G\).  Conjugation by \(H\) preserves
the characters \(\alpha,\delta\).  Conjugation by \(c\)
sends each to its inverse: \(R(c)=1\) and the disconnected
action is \(R(h)\mapsto J_3R(h)^{-t}J_3\), whose middle
entry and determinant are respectively
\(\alpha(h)^{-1},\delta(h)^{-1}\).
Triviality therefore transfers to every place above \(\ell\).

At a rational prime \(p\ne\ell\), restrict to
\(H\cap I_p\).  The cyclotomic character is unramified and
the residual model is the diagonal representation
\((\psi,1,\psi^{-1})\).  The finite character \(\psi\)
has order prime to \(\ell\), by \texttt{GD}.  The point's
\texttt{Kill} condition says exactly that the kernel of this
finite residual representation is killed by \(R\).
Thus \(R\) factors on that inertia through a finite group
of order prime to \(\ell\).  The same holds for \(\alpha\)
and \(\delta\).  A prime-to-\(\ell\) root \(b\) with
\(b\equiv1\) is \(1\): in
\(b^m-1=(b-1)(1+\cdots+b^{m-1})\) the second factor is a
unit if \(\ell\nmid m\).  Both characters are therefore
trivial on these inertia groups too.

For completeness, the residual-kernel description just used does
not lose characteristic-zero information.  Reduction is injective
on prime-to-\(\ell\) roots of unity by the preceding argument;
on \(H\cap I_p\) the kernel of the residual diagonal is
therefore the kernel of \(\psi\).  If desired, the whole
finite inertial action is conjugate to that diagonal action:
the averaging intertwiner
\(\sum_{q}R_0(q)^{-1}R(q)\) is congruent to \(|Q|1\)
and hence invertible.  Only the prime-to-\(\ell\) order,
not this stronger conjugacy, is needed above.

Class-field theory identifies the everywhere unramified abelian
quotient of \(G_K\) with the ideal class group, since \(K\)
has no real places.  Thus \(\alpha\) and \(\delta\)
factor through a group of order \(h_K\).  The assumption
\(\ell\nmid h_K\) and residual closeness again force them
to be trivial.  As \(\det R=\alpha\det B\), (KTR.10)
follows.  This uses the original \(K\), its class number, and
the original residual point, with no auxiliary coefficient-ring
or modularity hypothesis.
% END KTR5

% BEGIN KTR6
\paragraph{KTR6. The explicit original rank-two extension.}
\label{claim:KTR6}
There is a continuous representation \(W:G\to
\mathrm{GL}_2(\overline{\mathbf Q}_\ell)\) satisfying
the literal \(D.\texttt{IsExt}(r,W)\), and
\[
 W|_H=B,\qquad W(c)=D_2:=\operatorname{diag}(-1,1),
 \qquad \det W=\eta.
 \tag{KTR.11}
\]
For every \(g\in G\), \(R(g)\) is a block embedding of
an \(\mathrm{SL}_2\) matrix with middle entry \(1\).

\emph{Proof.}
Every \(g\notin H\) is \(hc\) with \(h\in H\).
The semidirect multiplication and \(R(c)=1\) give
\(R(hc)=R(h)\), so KTR5 proves the last assertion.
For \(g\in G\), let \(B_g\) be this corner matrix and put
\[
 W(g)=B_gD_2^{(1-\eta(g))/2}.
 \tag{KTR.12}
\]
Write \(J_2=\left(\begin{smallmatrix}0&1\\1&0\end{smallmatrix}\right)\).
For every \(B\in\mathrm{SL}_2\), direct multiplication gives
\(J_2B^{-t}J_2=D_2BD_2\).  The disconnected action in the
original \(\mathrm{GG}_3\) restricts on the corner to its
left-hand side.  Its semidirect multiplication therefore makes
(KTR.12) a homomorphism.  The two cosets of \(H\) are open;
on each coset its entries are fixed signed entries of \(R\).
This proves continuity.  Formula (KTR.11) follows immediately,
including \(\operatorname{tr}W(c)=0\).
The definition of \texttt{IsExt} asks exactly for the equality
with the corner on \(H\) and this trace at \(c\).
Thus this is an actual extension covered by the original universal
modularity hypothesis.  Its restriction to \(H\) is irreducible
by \texttt{TypeII}, hence \(W\) is irreducible too.
% END KTR6

% BEGIN KTR7
\paragraph{KTR7. An elementary two-dimensional self-twist criterion.}
\label{claim:KTR7}
Let \(V\) be an irreducible two-dimensional representation of
a group \(\Gamma\) over an algebraically closed field of
characteristic zero.  If \(N\lhd\Gamma\), \(\Gamma/N\)
is abelian, and \(V|_N\) is reducible, then there is a
nontrivial character \(\xi:\Gamma/N\to\{\pm1\}\) with
\(V\simeq V\xi\).

\emph{Proof.}
An \(N\)-stable line cannot be the unique such line, for
normality would make it \(\Gamma\)-stable.  Two distinct
stable lines give a diagonal decomposition
\(V|_N=\nu_1\oplus\nu_2\).
If \(\nu_1\ne\nu_2\), these are the only stable lines.
The action of \(\Gamma\) permutes them nontrivially, giving
\(\xi\); the matrix acting as \(1,-1\) on them intertwines
\(V\) and \(V\xi\).

If \(\nu_1=\nu_2\), all elements of \(N\) act by scalars.
Commutators of elements of \(\Gamma\) lie in \(N\), so
their matrices commute projectively.  Choose a nonscalar matrix
\(A\) in the image, possible by irreducibility.  For each
\(g\), conjugation sends \(A\) to \(\xi(g)A\).
Determinants give \(\xi(g)^2=1\), and conjugation shows that
\(\xi\) is a character trivial on \(N\).
If it were trivial, \(A\) would commute with the whole
irreducible representation and would be scalar, a contradiction.
Finally \(A\) itself is an invertible intertwiner for this
self-twist.  This proves the assertion without a semisimplicity
assumption on the initial restriction to \(N\).
% END KTR7

% BEGIN KTR8
\paragraph{KTR8. Irreducibility on the full infinite cyclotomic kernel.}
\label{claim:KTR8}
The actual \(W\) constructed in KTR6 is irreducible on
\(G_{K(\mu_{\ell^\infty})}\).

\emph{Proof.}
On the group \(J\) of KTR4, the first flag line has character
\(\epsilon^e\).  It lies in the corner plane: the middle
line is trivial by KTR5, whereas \(\epsilon(s_*)^e\ne1\).
The ordered flag is preserved by \(I\).  Its first line is
therefore an inertia-stable line for \(W\), with character
\(\kappa\), and \(\det W|_I=1\) gives the other diagonal
character \(\kappa^{-1}\).  In particular
\[
 \operatorname{tr}W(s)=\kappa(s)+\kappa(s)^{-1}\quad(s\in I),
 \qquad \kappa|_J=\epsilon^e.
 \tag{KTR.13}
\]
The character \(\kappa\) has infinite image, so choose
\(u\in I\) with \(\kappa(u)^8\ne1\).

Set \(N=\ker(\epsilon|_H)\).  If the restriction to \(N\)
were reducible, KTR7 would give a nontrivial quadratic self-twist
\(\xi\) trivial on \(N\).  The image of \(I\) under
\(\epsilon\) is all \(\mathbf Z_\ell^\times\), and
\(I\subset H\), so \(I\) surjects onto \(H/N\).
There is consequently \(s\in I\) with \(\xi(s)=-1\).
Equality of traces under the self-twist and (KTR.13) imply
\(\kappa(s)^4=1\).  The same reasoning applies to
\(su^2\), because \(\xi(u^2)=1\).  Dividing the two
fourth-power equalities gives \(\kappa(u)^8=1\), a
contradiction.  Thus the restriction to \(N\) is irreducible.
ADK2 identifies this actual closed subgroup with
\(G_{K(\mu_{\ell^\infty})}\); no finite cyclotomic level
has been substituted.  We have not needed to identify the finite
part of \(\kappa\) with the original \(\tau\) on all of
\(I\), nor to invoke a stronger local deformation theorem.
% END KTR8

\subsection{The corner adjoint and the nine-dimensional decomposition}

% BEGIN KTR9
\paragraph{KTR9. The actual modular extension has a cuspidal quadratic base change.}
\label{claim:KTR9}
Suppose the \(W\) in KTR6 satisfies the original
\(D.\texttt{IsModular}(\iota,W)\).  There is a regular
algebraic cuspidal representation \(\Pi\) of
\(\mathrm{GL}_2(\mathbf A_K)\), polarized over \(\mathbf Q\),
whose normalized Galois representation is the actual \(W|_H\).
At an unramified place \(w\) with norm \(q_w\), its
normalized Satake roots \(b_{w,i}\) satisfy
\[
       \operatorname{Eig}(W(F_w))
                  =\{\iota^{-1}(q_w^{1/2}b_{w,i})\}_{i=1,2}.
 \tag{KTR.14}
\]

\emph{Proof and precise published input.}
Modularity supplies an original newform \(f\) of weight
\(2e+1\) with attachment (KTR.2).  RNC2 constructs its actual
cuspidal adelization \(\pi^0\) over \(\mathbf Q\) and
the regular algebraic normalization
\(\pi=\pi^0|\det|^{1/2}\), whose normalized Galois
representation is \(W\).  Its infinitesimal parameters are
\(e+\tfrac12,\tfrac12-e\), distinct because \(e\ge1\).
This is a statement about the given \(f,W,\iota\), not an
existential replacement of its Frobenius system.

Apply quadratic base change for \(\mathrm{GL}_2\).
The precise primary input is Langlands, \emph{Base change for
GL(2)}, Annals of Mathematics Studies 96 (1980),
\S2(A)--(F), Lemma 11.3(b), and Proposition 11.4:
a cuspidal representation over a characteristic-zero number
field has a cuspidal lift to a given quadratic extension unless
it is induced from a character of that extension; the lift
matches local base change, including infinity and all spherical
places.  It suffices to apply the theorem to the unitary
\(\pi^0\), then apply the same norm twist.
The primary text is
\url{https://publications.ias.edu/sites/default/files/base-change-for-gl2_rpl_2.pdf}.
In the exceptional induced case, twisting by \(\eta\) fixes
the representation: the inducing character's induction is
unchanged because \(\eta|_H=1\).  This also follows directly
from its unramified parameters and strong multiplicity one.

That exceptional case cannot occur here.  If
\(\pi\simeq\pi\eta\), its unramified polynomials and
CAP6 give \(W\simeq W\eta\).  An intertwiner then commutes
with \(W(H)\), so is scalar by its irreducibility, while at
\(c\) it must anticommute with \(W(c)\).  An invertible
scalar cannot do so in characteristic zero.  Thus the base change
\(\Pi\) is cuspidal.  Its archimedean parameter is restriction
of the real parameter to \(\mathbf C^\times\); at each
complex embedding it has the same two distinct half-integral
parameters.  This is regular algebraicity in the original
Capelli/Harish--Chandra convention checked in ADK3.

For the normalization, let \(a_1,a_2\) be the spherical Satake
roots of \(\pi_p\).  At a place \(w\mid p\) of residue
degree \(f_w\), base change gives \(a_1^{f_w},a_2^{f_w}\).
Since \(q_w=p^{f_w}\),
\(q_w^{1/2}a_i^{f_w}=(p^{1/2}a_i)^{f_w}\).
This proves (KTR.14) for every unramified place outside one
finite set, including inert places.

Finally \(\det W|_H=1\) and \(W\) extends to \(G\).
Thus the normalized Galois roots at \(w^c\) equal those at
\(w\), and form an inverse pair.  Translating (KTR.14) back
to Satake roots gives
\[
 \{b_{w^c,i}\}=\{q_w^{-1}b_{w,i}^{-1}\}.
 \tag{KTR.15}
\]
Class-field theory supplies an algebraic Hecke character
\(\lambda_0\) over \(\mathbf Q\) whose pullback to \(K\)
has \(\ell\)-adic character \(\epsilon|_H\); it is the
norm character fixed by these geometric-Frobenius values.
If necessary multiply it by the quadratic Hecke character of
\(K/\mathbf Q\) so that its infinity value at \(-1\) is
\((-1)^2=1\).  This does not change its pullback or (KTR.15).
Strong multiplicity one now gives
\(\Pi^c\simeq\Pi^\vee\otimes(\lambda\circ N_{K/\mathbf Q}\circ\det)\).
It is exactly the polarized pair required by BLGGT \S2.1;
its normalized Galois multiplier on \(H\) is
\(\epsilon^{-1}\lambda_\ell=1\).
BLGGT Theorem 2.1.1 now constructs the normalized attached
representation.  Its polynomials agree with (KTR.14), so
Chebotarev and the semisimple character criterion identify it
with \(W|_H\).  Thus the pairing hypothesis is established
before invoking that existence theorem.
% END KTR9

% BEGIN KTR10
\paragraph{KTR10. Vanishing for the complete polarized corner adjoint.}
\label{claim:KTR10}
On \(M_2(\overline{\mathbf Q}_\ell)\), let \(A\) act by
\[
 A(h)X=W(h)XW(h)^{-1}\quad(h\in H),
 \qquad A(c)X=-J_2X^tJ_2.
 \tag{KTR.16}
\]
Then \(H^1_f(\mathbf Q,A)=0\) in the full original sense.

\emph{Proof.}
We verify the pairing rather than infer it from self-duality
alone.  Put
\(S=\left(\begin{smallmatrix}0&1\\-1&0\end{smallmatrix}\right)\).
Every two-by-two matrix satisfies
\(W(g)^t S W(g)=\det W(g)S=\eta(g)S\).
Since \(W(c)=D_2\) and \(S D_2=J_2\), it follows that
\[
             W(h)^tJ_2W(chc^{-1})=J_2\quad(h\in H).
 \tag{KTR.17}
\]
The pairing \(J_2\) is symmetric and nondegenerate.
KTR9 supplies the exact regular algebraic polarized cusp and
its attached representation.  BLGGT Theorem 2.1.1 supplies the
same normalized pairing; its ratio with (KTR.17) commutes with
the irreducible \(W(H)\), hence is scalar.  The disconnected
adjoint formula is therefore exactly (KTR.16), as in ADK4:
the differential of inverse transpose is minus transpose.
It acts on all of \(M_2\), including the scalar line.

KTR8 is precisely the full infinite-cyclotomic irreducibility
hypothesis of Thorne's Theorem 6.1, with \(F=K\),
\(F^+=\mathbf Q\), and \(n=2\).
That theorem states \(H^1_f(\mathbf Q,A)=0\) for this
polarized adjoint.  Its scope has no ordinary hypothesis and
no exclusion of the prime \(3\).  The exact published
statement, its full \(M_n\) action, and its BLGGT conventions
were checked in ADK3--ADK6; here (KTR.14)--(KTR.17) and KTR8
verify every representation-dependent hypothesis afresh.
KTR1 brings the result from finite coefficients and the classical
period ring to every cocycle and every local primitive in the
original \texttt{Af.IsBK} definition.  Thus the conclusion
is actual coboundary vanishing, not merely a vanishing dimension
for a different coefficient realization.
% END KTR10

% BEGIN KTR11
\paragraph{KTR11. The remaining quadratic scalar line.}
\label{claim:KTR11}
The one-dimensional representation with character \(\eta\)
has \(H^1_f(\mathbf Q,\overline{\mathbf Q}_\ell(\eta))=0\).

\emph{Proof.}
By KTR1 work over a finite coefficient field \(E\).
Restriction to \(K\) changes the representation into the
trivial representation.  For a trivial representation the local
finite condition at a place above \(\ell\) is the unramified
condition.  Explicitly the Bloch--Kato fundamental exact sequence
used in BX7 gives \(\dim_E H^1_f(K_v,E)=1\): the de Rham
tangent quotient is zero and the invariant space has dimension
one.  The unramified group is already one-dimensional and is
contained in the finite group, so they agree.  Away from
\(\ell\) this equality is the definition.  A global cocycle
for the trivial representation is a continuous additive
homomorphism.  If it is finite everywhere it therefore kills
every inertia subgroup, and factors through the finite ideal
class group of \(K\).  There is no nonzero homomorphism from
that finite group to the additive group of a characteristic-zero
field.  Thus \(H^1_f(K,E)=0\).  Restriction from \(\mathbf Q\)
is injective by the finite-group averaging argument in KTR3,
which proves the assertion.  This rational vanishing itself
does not require the stronger assumption \(\ell\nmid h_K\).
% END KTR11

% BEGIN KTR12
\paragraph{KTR12. Explicit equivariant decomposition of the nine-dimensional adjoint.}
\label{claim:KTR12}
The full original \(\operatorname{ad}r\) is the direct sum
of \(A\), \(W\), \(W\eta\), and \(\eta\), through
the following concrete continuous linear maps.  For
\(X\in M_3\), \(Y\in M_2\), and \(v=(v_0,v_1)\), set
\[
 \begin{aligned}
 p_A X&=\begin{pmatrix}X_{00}&X_{02}\\X_{20}&X_{22}\end{pmatrix},
 &s_A Y&=\begin{pmatrix}Y_{00}&0&Y_{01}\\0&0&0\\Y_{10}&0&Y_{11}\end{pmatrix},\\
 p_tX&=(X_{01}+tX_{12},\ X_{21}-tX_{10}),
 &s_tv&=\begin{pmatrix}0&v_0&0\\-t v_1&0&t v_0\\0&v_1&0\end{pmatrix},
 \quad t\in\{1,-1\},\\
 p_CX&=X_{11},&s_Cz&=zE_{11}.
 \end{aligned}
 \tag{KTR.18}
\]
Here \(p_1,s_1\) intertwine with \(W\), and
\(p_{-1},s_{-1}\) with \(W\eta\).  The identity is
\[
       X=s_Ap_AX+\tfrac12s_1p_1X+
                       \tfrac12s_{-1}p_{-1}X+s_Cp_CX.
 \tag{KTR.19}
\]

\emph{Proof.}
On \(H\), write the corner matrix as
\(B=\left(\begin{smallmatrix}a&b\\c&d\end{smallmatrix}\right)\)
with \(ad-bc=1\), and its three-dimensional embedding as
\(\widetilde B\).  The original adjoint action there is
\(X\mapsto\widetilde B X\widetilde B^{-1}\).
The four entries of \(p_A\) undergo \(Y\mapsto BYB^{-1}\).
For the off-diagonal entries, direct multiplication gives
\[
 \begin{split}
 p_t(\widetilde B X\widetilde B^{-1})_0
       &=a(X_{01}+tX_{12})+b(X_{21}-tX_{10}),\\
 p_t(\widetilde B X\widetilde B^{-1})_1
       &=c(X_{01}+tX_{12})+d(X_{21}-tX_{10}).
 \end{split}
\]
The same multiplication gives
\(\widetilde B s_t(v)\widetilde B^{-1}=s_t(Bv)\)
and fixes \(X_{11}\).

At \(c\), the source adjoint is
\(\Theta(X)=-J_3X^tJ_3\), with \(J_3\) the reverse
identity matrix; this follows from \(R(c)=1\) and
\(\eta(c)=-1\).  Its corner action is
\(-J_2(p_AX)^tJ_2\), and its middle entry is \(-X_{11}\).
For the other two maps one obtains
\[
       p_t\Theta(X)=\operatorname{diag}(-t,t)\,p_tX.
 \tag{KTR.20}
\]
For \(t=1\) this is \(D_2=W(c)\); for \(t=-1\) it
is \(-D_2=(W\eta)(c)\).  Substitution in \(s_t\)
gives the identical intertwining law in the other direction.
The groups \(H\) and \(\{c\}\) generate \(G\), so
all these maps intertwine the full disconnected action.
Finally (KTR.19) is an entrywise identity: the sum of the
\(t=1\) and \(t=-1\) terms cancels each unwanted cross
entry and retains the desired one.  Their pairwise projections
have zero cross compositions and \(p_ts_t=2\,1\), while
\(p_As_A=p_Cs_C=1\).  This proves the direct-sum assertion.
All maps have fixed rational coefficients and are continuous
for the original product \(\ell\)-adic topology.
% END KTR12

% BEGIN KTR13
\paragraph{KTR13. The full literal RigII producer.}
\label{claim:KTR13}
For every original datum and actual point of \texttt{RigII},
if \(\texttt{TypeII}(\eta,r)\) and
\[
 \forall W,\quad D.\texttt{IsExt}(r,W)
                 \ \Longrightarrow\ D.\texttt{IsModular}(\iota,W),
 \tag{KTR.21}
\]
then the source's
\(\texttt{BlochKatoSelmerVanishes}(\ell,\operatorname{ad}r)\)
holds.

\emph{Proof.}
KTR4--KTR6 construct the actual \(W\) from the point and
verify \texttt{IsExt}; (KTR.21) gives its modular newform for
the original fixed \(\iota\).  KTR2--KTR3 give Selmer
vanishing for \(W\) and \(W\eta\).  KTR8--KTR10 give
vanishing for the corner adjoint \(A\); KTR11 gives it for
the remaining scalar character \(\eta\).

Let \(a:G\to M_3(\overline{\mathbf Q}_\ell)\) be any
continuous cocycle satisfying the original local conditions.
Apply \(p_A,p_1,p_{-1},p_C\) from KTR12.  KTR1 makes all
four images continuous Selmer cocycles for their stated actions.
Choose actual primitives \(Y,v_+,v_-,z\) supplied by the
four vanishings.  Equivariance and (KTR.19) show that
\[
 a(g)=(\operatorname{ad}r)(g)Q-Q,
 \qquad
 Q=s_A Y+\tfrac12s_1v_++\tfrac12s_{-1}v_-+s_Cz.
 \tag{KTR.22}
\]
This is the literal required global coboundary.

The argument retains the supplied \(\Omega,K,\ell,\mathcal O\),
the complete \(\texttt{GD}\) conditions, \(N\ne0\),
\(e\), \(\psi\), the universal local ring and its actual
point map, the marked \(g_1\), and the normalization at the
specified \(c\).  It uses only the subset of those hypotheses
identified in KTR4--KTR8, so it imposes no new level or weight
condition.  It retains the full infinite cyclotomic kernel,
the actual representation, every local \texttt{IsBK} witness,
and all nine adjoint coordinates.  The original odd-prime
scope includes \(3\); no division by \(\ell\) or omitted
small-prime case is introduced.  The ordinary-coefficient
automorphy producer and the remaining \texttt{OneChar} length
argument are distinct tasks and are not conclusions of this
claim.
% END KTR13

\subsection{Source and foundation boundary}

The original endpoint statements are
\path{ModularityArgument/KatoQuadraticCharacters/KatoTwoOfUseN.lean},
lines 11--53.  The local ideal and its ordered products are in
\path{PadicHodgeTheory/IsInertiaInfiniteCyclotomic.lean},
lines 87--138.  The actual point and extension predicates are
\path{UnitaryThreeGroups/RhoKApply.lean}, lines 16--93.
The finite-dimensional source route appears in
\path{ModularityArgument/KatoQuadraticCharacters/ReducibleThreeDim.lean};
we have written its actual projection, section, and primitive
calculations in KTR12--KTR13.  KTR4--KTR8 give a smaller local
argument sufficient for this endpoint, without claiming the
stronger full local crystalline-component theorem.

The new published applications are Kato Theorem 14.2(1), with
its \S14.1 local conditions and the actual critical twist, and
quadratic \(\mathrm{GL}_2\) base change in Langlands's
Lemma 11.3(b) and Proposition 11.4.  BLGGT Theorem 2.1.1 and
Thorne Theorem 6.1 are the previously matched foundational
theorems; KTR9--KTR10 check this new representation's exact
hypotheses.  The finite-coefficient and tied-period comparisons
are inherited from ADK1--ADK5, BX1, BX7 and DX1; the actual
newform adelization is RNC2 and the dense Frobenius comparison
is CAP6.  All these dependencies and every exact source range
are recorded separately from the proof in the component
provenance.  No original Lean proof is used as the human proof.
